\textbf{Pour cette première partie, aucune justification n'est demandée et une seule réponse est possible par question. Pour chaque question, reportez son numéro sur votre copie et indiquez votre réponse.}

Une réponse fausse ou l'absence de réponse n'enlève aucun point.

\subsection*{Question 1}

Déterminer la médiane de la série : $12$ ; $9$ ; $7$ ; $23$ ; $9$ ; $25$ ; $7$.



\subsection*{Question 2}

\medskip

Donner la notation scientifique de $\np{0,000457}$.



\begin{minipage}{0.48\linewidth}
\subsection*{Question 3}

\medskip

Calculer l'aire, en cm$^2$, du triangle ci-contre.

\end{minipage}\hfill
\begin{minipage}{0.48\linewidth}
\begin{center}
%Début Tikz
%\begin{tikzpicture}[scale=1]
%% Vertices: A (left), B (top right), C (bottom right)
%\coordinate (A) at (0,0);
%\coordinate (B) at (3,1.5);
%\coordinate (C) at (3,-1.5);
%
%\draw (A) -- (B) -- (C) -- cycle;
%
%% Right angle at B
%\draw (B) +(-0.2,-0.1) -- +(-0.15,-0.28) -- +(-0.05,-0.2);
%
%% Labels
%\node[left] at (A) {$A$};
%\node[above right] at (B) {$B$};
%\node[below right] at (C) {$C$};
%
%% Side labels
%\node[above left] at ($(A)!0.5!(B)$) {$8\ \text{cm}$};
%\node[right] at ($(B)!0.5!(C)$) {$6\ \text{cm}$};
%\node[below] at ($(A)!0.5!(C)$) {$10\ \text{cm}$};
%\end{tikzpicture}
%fin Tikz
%début PSTricks
\psset{unit=1cm}
\begin{pspicture}(5.8,3.6)
\uput[l](0.4,1.4){A}\uput[dr](5.4,0.4){C}\uput[u](4.34,3.){B}
%\psarc(2.9,0.9){2.5495}{0}{90}
\pspolygon(0.4,1.4)(4.34,3.)(5.4,0.4)%ABC
\rput{-157}(4.34,3.){\psframe(0.25,0.25)}
\rput{23}(2.6,2.6){8 cm}\rput(5.3,1.9){6 cm}\rput{-11}(2.9,0.7){10 cm}
\end{pspicture}
% fin PSTricks
\end{center}
\end{minipage}





\subsection*{Question 4}

\medskip

Une boîte opaque contient des beignets tous identiques, garnis de confitures différentes :
\begin{itemize}[label=$\bullet~$]
\item 6 beignets sont à l'abricot ;
\item 5 beignets sont à la pomme ;
\item 4 beignets sont à la framboise.
\end{itemize}

Déterminer la probabilité de piocher au hasard un beignet à la framboise.



\subsection*{Question 5}

\medskip

Un article coûte $800$~\euro{}. Son prix baisse de $10~\,\%$.

Calculer le prix, en euro, de l'article après réduction.



\subsection*{Question 6}

\medskip

Développer et réduire l'expression $B = 4y(3y - 1)$.



\subsection*{Question 7}

\medskip

Recopier la réponse permettant de compléter l'égalité 3,57 L = \ldots cm$^3$
.


\subsection*{Question 8}

\medskip

Recopier sur la copie l'image de 4 par la fonction affine $f$ définie par $f(x) = 3x - 5$.

\begin{tblr}{colspec = {*4{X[c]}},hlines,vlines}
	\textbf{A.}\quad 3 &
	\textbf{B.}\quad 7 &
	\textbf{C.}\quad 12 &
	\textbf{D.}\quad 29\\
\end{tblr}

\begin{minipage}{0.65\linewidth}
\subsection*{Question 9}

Le quadrilatère ABCD ci-contre est tracé à main levée.

À partir des codages donnés, en déduire sa nature parmi
les quatre réponses proposées et la recopier sur la copie.
\end{minipage}
\hfill
\begin{minipage}{5cm}
\psset{unit=0.85cm}
\begin{pspicture}(0,-0.5) (3.,4.5)
%\psgrid
\uput[dl](0.2,0.2){A}\uput[dr](4,0.8){B}\uput[ur](3.8,4){C}\uput[ul](0,3.4){D}
\pslineByHand(0.2,0.2)(4,0.8)(3.8,4)(0,3.4)(0.2,0.2)
\pslineByHand(0.2,0.2)(3.8,4)
\pslineByHand(4,0.8)(0,3.4)
\def\barre{\psline(-0.05,-.2)(-0.05,.1)}
\def\barbarre{\psline(-0.05,-.2)(-0.05,.2)\psline(0.05,-.2)(0.05,.2)}
\rput{35}(1.2,1.3){\barre}\rput{35}(2.9,3.1){\barre}
\rput{-60}(1.3,2.5){\barbarre}\rput{-60}(3,1.45){\barbarre}
\end{pspicture}
\end{minipage}

\begin{tblr}{colspec = {*2{X[c]}},hlines,vlines}
	\textbf{A.}\quad Un losange. &
	\textbf{B.}\quad Un rectangle. \\
	\textbf{C.}\quad Un carré. &
	\textbf{D.}\quad Un parallélogramme. \\
\end{tblr}
